\section{Introduksjon (10 poeng)}\label{introduksjon-10-poeng}
Undergangen av Titanic i 1912 er en av verdens mest kjente maritime tragedier. Dette prosjekt tar for seg datasettet \textit{Titanic-Dataset.csv}, hentet fra kaggle \parencite[]{titanic_dataset} som dokumenterer informasjon knyttet til 891 passasjerer ombord på skipet.

Prosjektet vårt er delt i to deler. Den første delen benytter multippel lineær regresjon for å undersøke om alder kan estimeres basert på passasjerklasse \textit{(Pclass)}, familiestørrelse \textit{(FamilySize)}, og kjønn \textit{(Sex)}. 
% Bakgrunnen for dette er valgt er at en stor andel av radene mangler registrert alder (rundt ~20\%). 
Den andre delen tar for seg KNN klassifisering for å undersøke hvor godt overlevese kan klassifiseres basert på kjønn \textit{(Sex)}, alder \textit{(Age)}, passasjerklasse \textit{(Pclass)} og familiestørrelse \textit{(FamilySize)}. 
 
De to delene til prosjektet er som følger:
\begin{itemize}
    \item \textbf{Multippel lineær regresjon (alder)}: Datasettet mangler verdier knyttet til enkelte variabler som alder (\textit{Age}). Ettersom at vi ønsker å benytte alder som en forklaringsvariabel ved klassifisering, estimerer vi alder der den mangler slik at vi kan få et bedre datagrunnlag. Vi ønsker derfor å undersøke i hvilken grad multippel lineær regresjon kan brukes til å estimere alder basert på passasjerklasse, familiestørrelse og kjønn.
    \item \textbf{Klassifisering, k-nærmeste-nabo}: Vi ønsker å se hvor godt passasjerenes overlevelse (\textit{Survived}) kan klassifiseres ved hjelp av KNN. Det vil også undersøkes hvordan valg av forklaringsvariabler og håndtering av manglende aldersverdier kan påvirke resultatene.
\end{itemize}

Formålet med prosjektet er å oppnå en bredere forståelse for hvordan sosiodemografiske faktorer kan forklare ulikhetene i overlevelse under en krisesituasjon. Gjennom prosjektet ønsker vi å finne ut: I hvor stor grad kan multippel lineær regresjon forklare/estimere alderen til personer med hjelp av flere uavhengige variabler som beskrevet/nevnt over og hvilke faktorer utgjør størst forskjell for sannsynligheten for overlevelse.

Målgruppen for prosjektet er individer som ønsker praktisk og forståelig innsikt i dataanalyse og klassifisering. I stedet for å finne modellen med størst mulig treffprosent, så ønsker vi å finne en metode som leverer gode resultater samtidig som den er enkel å forklare og tolke.